\section{问题重述}

\subsection{问题背景}

大语言模型的性能与模型规模、训练数据及计算资源密切相关。随着模型规模和训练数据量不断增加，训练所需的计算成本也随之上升。已有研究通过标度律描述参数规模、数据量和计算投入与训练损失之间的关系，为模型训练中的资源分配提供了依据。然而，实际训练数据在质量和领域组成上存在差异，即使采用相近的参数规模和计算量，不同的数据选择也可能产生不同的训练效果。

在算力有限的情况下，扩大模型规模、增加训练数据、提高数据质量和延长上下文长度都需要额外的计算投入，因而必须权衡不同资源配置带来的收益。此外，训练损失的降低并不能完全反映模型在推理、数学和代码等下游任务中的能力变化，仅根据损失指标进行资源配置，难以全面评价模型的实际性能。

本题围绕有限算力下的大语言模型能力提升问题，结合训练语料、数据配比实验、模型训练记录及能力评测数据，研究数据质量和领域组成对训练效果的影响，建立多因素标度关系，确定不同算力预算下的资源配置方案，并分析模型能力的历史演进及未来发展趋势。

\subsection{问题提出}

\textbf{问题一：}
利用附件 A 提供的训练语料质量指标，对不同来源的文本进行质量评价，分析各指标之间的评价冲突，并检验评价结果在抽样数据和扩展数据上的一致性。在此基础上，结合不同领域的训练数据配比实验，研究数据质量、领域配比及领域组合对训练损失的影响，建立相应的量化关系，并考察其在给定检验数据上的适用性。

\textbf{问题二：}
根据问题一的数据质量评价与领域配比分析结果，结合附件 B 中的模型训练记录，建立同时考虑参数规模、训练数据量、数据质量及领域配比的广义标度关系。通过不同训练轨迹和模型数据检验模型的有效性，分析各因素对训练损失的影响，以及数据质量提升与模型规模扩张之间的替代关系，并讨论模型向更大规模外推时的适用范围。

\textbf{问题三：}
基于前两问建立的模型，在不同算力预算下确定模型参数量、训练数据量、数据质量和领域配比的合理配置，使预测训练损失尽可能降低。计算成本需要同时考虑基础训练、数据质量提升和长上下文带来的额外开销。在此基础上，分析预算及质量成本变化对配置结果的影响，研究资源配置随预算增长发生的结构性变化，并确定注意力计算开销与基础训练开销相当时的临界上下文长度。

\textbf{问题四：}
利用附件 C 中的模型排行榜、历史记录及能力评测数据，研究开源大语言模型的能力演进，分析模型规模扩张与非规模技术进步对能力提升的贡献。在算力增长放缓的情景下，预测未来 12 个月或 24 个月的模型能力前沿，并评估预测结果的不确定性。同时，研究训练损失与下游任务评测得分之间的关系，分析二者的映射误差对能力演进分析和前沿预测的影响。
